% My CHEAT SHEET TO VIM
% Author: Bruno Kemmer
% Last Edit: 2025-01-29
%
% DISCLAIMER: This work is licensed under a Creative Commons Attribution-NonCommercial-ShareAlike 3.0 Unported License (http://creativecommons.org/licenses/by-nc-sa/3.0/). The cheatsheet template was adapted from Winston Chang's version found at http://www.stdout.org/~winston/latex/. This cheatsheet is intended for non-commercial use.


\documentclass[10pt,landscape]{article}
\usepackage{multicol}
\usepackage{ifthen}
\usepackage[landscape]{geometry}
\pagestyle{empty}
\setcounter{secnumdepth}{0}
\setlength{\parindent}{0pt}
\setlength{\parskip}{0pt plus 0.5ex}


% This sets page margins to .5 inch if using letter paper, and to 1cm       
% if using A4 paper. (This probably isn't strictly necessary.)              
% If using another size paper, use default 1cm margins.
\ifthenelse{\lengthtest { \paperwidth = 11in}}
	{ \geometry{top=.5in,left=.2in,right=.2in,bottom=.5in} }
	% {\ifthenelse{ \lengthtest{ \paperwidth = 297mm}}
		% {\geometry{top=1cm,left=1cm,right=1cm,bottom=1cm} }
		% {\geometry{top=1cm,left=1cm,right=1cm,bottom=1cm} }
	% }
 
% Redefine section commands to use less space 
\makeatletter
\renewcommand{\section}{\@startsection{section}{1}{0mm}%
                                {-1ex plus -.5ex minus -.2ex}%
                                {0.5ex plus .2ex}%x
                                {\normalfont\large\bfseries}}
\renewcommand{\subsection}{\@startsection{subsection}{2}{0mm}%
                                {-1explus -.5ex minus -.2ex}%
                                {0.5ex plus .2ex}%
                                {\normalfont\normalsize\bfseries}}
\renewcommand{\subsubsection}{\@startsection{subsubsection}{3}{0mm}%
                                {-1ex plus -.5ex minus -.2ex}%
                                {1ex plus .2ex}%
                                {\normalfont\small\bfseries}}
\makeatother

%%%%%%%%%%%%%%%%%%%%%%%%%%%%%%%%%%%%%%%%%%%%%%%%%%%%%%%%%%%%%%%

\begin{document}
\raggedright
\begin{multicols}{3}

% multicol parameters                                                       
% These lengths are set only within the two main columns                    
\setlength{\columnseprule}{0.25pt}
\setlength{\premulticols}{1pt}
\setlength{\postmulticols}{1pt}
\setlength{\multicolsep}{1pt}
\setlength{\columnsep}{2pt}

\begin{center}
     \Large{\underline{\textbf{My VIM Cheat Sheet}}} \\
\end{center}


\verb! !\\

\section{Plugins}
\verb! !\\
\subsection{NERDTree}

\verb!<Leader> t! \hfill Focus to NERDTree window\\
\verb!<Leader> e! \hfill Toogle NERDTree window\\
\verb!<C-f>! \hfill Highlight the current file on NERDTree window\\


\verb! !\\

\subsection{Jedi-vim}

\verb!<C-Space>! \hfill Python Auto-Completion\\
\verb!<C-n> / <C-p>! \hfill Next/Previous suggestions\\
\verb!<Leader>g! \hfill Goto assignment\\
\verb!<Leader>d! \hfill Goto definition\\
\verb!<Leader>s! \hfill Goto (typing) stub\\
\verb!K! \hfill Show Documentation/Pydoc\\
\verb!<Leader>r! \hfill Renaming\\
\verb!<Leader>n! \hfill Show usages\\
\verb!:Pyimport os! \hfill Opens the `os` module\\

\verb! !\\
\subsection{Vim-Commentary}

\verb!gcc! \hfill Comment a line\\
\verb!gc<motion>! \hfill Comment the target motion\\
\verb!gc! \hfill On visual: Comment the selection\\

\verb! !\\
\subsection{Vim-slime}
\verb!<C-c><C-c>! \hfill Send text to REPL\\
\verb!<C-c>v! \hfill Reconfigure target\\



% \subsection{Vscode Debug / Run}

% \verb!<C-Shift-d>! \hfill Show debug tab.\\
% \verb!<C-Shift-Y>! \hfill Toggle debug console.\\
% \verb!<leader> r! \hfill Run python code on terminal.\\
% \verb!ge! \hfill On visual mode: Evaluate on debug console.\\
% \verb!gaw! \hfill On visual mode: Add selection to watch list.\\
% \verb!gw / <C-1>!\hfill Focus on watch list tab / goes back.\\
% \verb!! \hfill \\
% \begin{multicols}{2}
%  \verb!F5! \hfill Start debugging.\\
%  \verb!F5! \hfill Continue debugging.\\
%  \verb!F10! \hfill Step over.\\
%  \verb!F11! \hfill Step into.\\
%  \verb! Shift-F11! \hfill Step out.\\
%  \verb! Shift-F5! \hfill Stop.\\
%  \verb! F9! \hfill Add a breakpoint to the current line.\\
% \end{multicols}
% \columnbreak

% \subsection{Vscode bookmarks extension}
% \verb!<leader> b! \hfill List bookmarks.\\
% \verb!<leader> m! \hfill Toggle bookmark.\\
% \verb!<leader> l! \hfill Toggle labeled bookmark.\\
% \verb!<leader> n/N! \hfill Next/Previous bookmark.\\
% \verb!! \hfill \\
% \subsection{Vscode Github extension}
% \verb!<C-Shift G>! \hfill Open github tab.\\
% \verb!<leader> g a! \hfill Open actions page.\\
% \verb!<leader> g r! \hfill Open repository page.\\
% \verb!<leader> g c! \hfill Open commits page.\\
% \verb!<leader> g y! \hfill Copy a permalink of the current line.\\
% \verb!! \hfill \\

% \textbf{Remember: }\textit{These can all be used with a repeater.}
\verb! !\\
\subsection{VIM}

\verb!<leader>! \hfill \textbf{space}\\
\verb!.! \hfill Repeats previous command\\

\verb! !\\
\subsection{Movement}

\verb!<C-d>! \hfill Move half-screen down.\\
\verb!<C-u>! \hfill Move half-screen up.\\
\verb!<C-f>! \hfill Move full screen forward.\\
\verb!<C-b>! \hfill Move full screen backward.\\
\verb!H / M / L! \hfill Bring cursor to top/middle/bottom of screen.\\
\verb!3H! \hfill Bring cursor to 3 lines from top of screen.\\
\verb!3L! \hfill Bring cursor to 3 lines from bottom of screen\\
\verb!}! \hfill Go forward by paragraph (the next blank line)\\
\verb!{! \hfill Go backward by paragraph (the previous blank line)\\

\verb! !\\
\verb!<C-w> + h/j/k/l or <C-number>! \hfill Move to * window.\\
\verb!<C-w> + H/J/K/L! \hfill Move window to * side.\\                                   
\verb!<C-w> + +/-! \hfill Increase/Decrease window height by 1.\\      
\verb!<C-w> -! \hfill Decrease window height by 1.\\                        
\verb!<C-w> + >/<!  \hfill Increase/Decrease window width by 1.\\                      
\verb!<C-w> + >!  \hfill Increase window width by 1.\\     
\verb!Ctrl + w + <!  \hfill Decrease window width by 1.\\           
\verb!<C-w> + =!  \hfill Equalize height and width of all windows.\\

\verb! !\\
\subsection{Buffers, Windows and Tabs}

\verb!:ls! \hfill List all buffers in session.\\
\verb!:bn/:bp! \hfill Navigate to next/previous buffer.\\
\verb!:bp! or \verb!:b#! \hfill Navigate to previous buffer.\\
\verb!:bf/:bl! \hfill Navigate to first/last buffer.\\
\verb!:bd! \hfill Delete current buffer.\\
\verb!:bd1! \hfill Delete buffer labeled "1".\\ 
\verb! !\\

\verb!gt! \hfill Next tab\\
\verb!gT! \hfill Previous tab\\
\verb!<C-t>! \hfill :tabnew\\
\verb!<C-^>! \hfill Return to previous file\\
\verb!<Leader>wo! \hfill Open current split in new tab\\
\verb!<C-hjkl>! \hfill Switch between splits\\


\verb! !\\


\verb!:vs ~/file! \hfill Open file in vertical pane.\\
\verb!:sp ~/file! \hfill Open file in horizontal pane.\\   
\verb!:vert sb2! \hfill Open buffer 2 in vertical pane.\\
\verb!:sb2! \hfill Opens buffer 2 in horizontal pane.\\
\verb! !\\

\verb!:e ~/file! \hfill Edit a file.\\
\verb!:tabe ~/file! \hfill Edit file in new tab.\\

\verb! !\\
\subsection{Editing}
\verb!J! \hfill Join lines with comments\\
\verb!w/b! \hfill Go to next/previous delimiter.\\
\verb!W/B! \hfill Go to next/previous word.\\
\verb! !\\
\verb!$! \hfill Go to end of line.\\
\verb!^! \hfill Go to first letter of a line.\\
\verb!0! \hfill Go to beginning of a line.\\
\verb!gg/G! \hfill Go to beginning/end of the file.\\
\verb!{/}! \hfill Go to beginning/end of paragraph.\\ 
\verb!:8! \hfill Go to line 8.\\
\verb! !\\
\verb!fP/FP! \hfill Find next/previous P in the line.\\
\verb!tP/TP! \hfill Find delimiter before/after next P.\\
\verb!;/,! \hfill Go to next/previous find.\\
\verb! !\\
\verb!yiw! \hfill Yank(copy) inside word\\
\verb!viwp! \hfill Select inside word and paste\\
\verb!viw"0p! \hfill Select inside word and paste content of register 0\\
\verb!di"/ci"! \hfill Delete/change text inside "".\\
\verb!da"/ca"! \hfill Delete/change text inside "", delete "".\\
\verb!u! \hfill Undo change.\\
\verb!Ctrl + r! \hfill Redo change.\\
\verb!v! \hfill Visual mode. \verb!j/k! to highlight text.\\
\verb!V! \hfill Visual Line mode. Highlight full lines.\\
\verb! !\\
\textit{In Visual mode, pressing \% under a bracket highlights until matching bracket.}\\
\verb!>>/<<! \hfill Indent/unindent in normal mode.\\
\verb!=3j! \hfill Autoindent 3 lines after cursor.\\

% \columnbreak


\subsection{Search and Replace}
\verb!:s/foo/bar/g! \hfill Replace each foo by bar (current line).\\
\verb!:%s/foo/bar/g! \hfill Replace each foo by bar (all lines).\\
\verb!:%s/foo/bar/gc! \hfill Replace each foo by bar (asking).\\
\verb! ! \\
\verb!*! \hfill Search for word under the cursor.\\
\verb!/[searchterm]! \hfill Search. \verb!n/N! for next/previous.\\
\verb!?[searchterm]! \hfill Backward search.\\
\textit{You can use regular expressions and escape chars.}\\
\textit{Can be used as motions with delete, change, copy.}\\
\verb!/! \textit{can be replaced by other separators,} \verb!%! \verb!#! etc.\\

% \verb! !\\
\subsection{Marks}

Marks allow you to jump to designated points in your code.
\verb!m{a-z}! \hfill Set mark {a-z} at cursor position\\
\verb!m{A-Z}! \hfill Sets a global mark (across files)\\
\verb!'{a-z}! \hfill Goto start of the line of the mark\\
\verb!''! \hfill Go back to the previous jump location\\


\columnbreak

% \verb! !\\
\subsection{Macros and Registers}
\verb!q<register><commands>q! \hfill Record a macro in register\\
\verb!qr! \hfill Start recording a macro in register "r" (q quits).\\
\verb!@r! \hfill Execute macro in register "r".\\
\verb!@@! \hfill Execute previous macro.\\
\verb! !\\

\verb!:reg! \hfill List all registers.\\
\verb!"r! \hfill Access content of register r.\\
\verb!"ry! \hfill Yank to register r.\\
\verb!""! \hfill Unnamed register has the content of the last yank or deleted text.\\
\verb!"0 to "9! \hfill Previous yanked of deleted text\\
\verb!".! \hfill The last inserted text\\
\verb!":! \hfill The last command executed (@: runs again)\\
\verb!"+ and "*! \hfill Have the systems clipboard\\
\verb!<C-r><register>! \hfill On insert mode: paste value of reg\\





\verb! !\\
\subsection{Miscellaneous}
\verb!:se ft=python! \hfill sets filetype.\\
\verb!:so ~/file! \hfill rebuilds a file.\\
\verb!:so %! \hfill rebuilds this file.\\
% not implemented on vscode yet!
\verb!:sav ~/file! \hfill Saves file elsewhere.\\
\verb!:!!\verb!ls -lah! \hfill Access the command line.\\
\verb!:r !!\verb!date! \hfill Append command output to cursor.\\
\verb! !\\
\verb=:%!jq . --tab= \hfill Prettify JSON file with tabs using jq.\\


\verb! !\\
\subsection{Documentation}
\verb!<C-]! \hfill Follow documentation hyperlink.\\
\verb!<C-o>! \hfill Go back.\\

% \textbf{To comment a block of text:}
% Go to the start of the first line, and in Visual block mode: \verb!<C-V>! select all lines, \verb!<S-I>! insert the comment char in the fist line and \verb!<CR>! after a second it will appear on the other lines.

\verb! !\\
\subsection{Folds}
\verb!zc! \hfill Close Fold.\\
\verb!zo! \hfill Open Fold.\\
\verb!za! \hfill Toogle Fold.\\
% \columnbreak

\verb! !\\
\subsection{Equivalents}
\begin{center}
\begin{tabular}{|r|l|}

\hline
\verb!c! & \verb!c$!          \\
\verb!s! & \verb!cl!          \\
\verb!S! & \verb!^c!          \\
\verb!I! & \verb!^i!          \\
\verb!A! & \verb!$a!          \\
\verb!o! & \verb!A<CR>!       \\
\verb!O! & \verb!ko!          \\
\hline
\end{tabular}
\label{table:ta}
\end{center}


\columnbreak


% \verb! !\\
% \verb! !\\
\section{Tmux}
\verb! !\\
\subsection{Sessions}
\verb!default prefix changed to <C-a>! \\
\verb!tmux! \hfill starts a new session.\\
\verb!tmux new -s NAME! \hfill starts it with that name.\\
\verb!tmux ls! \hfill lists the current sessions.\\
\verb!<C-a> d! \hfill detaches the current session.\\
\verb!tmux a! \hfill attach to last session.\\
\verb! !\\

\subsection{Windows}
\verb!<C-a> c! \hfill Creates a new window. To close it you can just terminate the shells doing \verb!<C-d>!.\\
\verb!<C-a> N! \hfill Go to the N th window. Note they are numbered.\\
\verb!<C-a> p! \hfill Goes to the previous window.\\
\verb!<C-a> n! \hfill Goes to the next window.\\
\verb!<C-a> ,! \hfill Rename the current window.\\
\verb!<C-a> w! \hfill List current windows.\\
\verb! !\\

\subsection{Panes}
\verb!<C-a> |! \hfill Split the current pane horizontally. Default \textit{"}.\\
\verb!<C-a> -! \hfill Split the current pane vertically. Default \textit{\%}.\\
\verb!<M-direction>! \hfill Move to the pane to direction.\\
\verb!<C-a> z! \hfill Toggle zoom for the current pane. \\
\verb!<C-a> [! \hfill Start scrollback. You can then press \verb!<space>! to start a selection and \verb!<enter>! to copy that selection. \\
\verb!<C-a> <space>! \hfill Cycle through pane arrangements. \\
\verb! !\\
\subsection{Multiple Monitors}
Create sessions for multiple monitors:
\verb!tmux new -s left_monitor!\\
\verb!tmux new -r left_monitor -s right_monitor!\\
\verb! !\\
\verb! !\\ Next item
\verb! !\\
\verb! !\\
\verb! !\\
\verb! !\\
\verb! !\\
\verb! !\\
\verb! !\\




% column 1
\columnbreak
\section{}
\verb! !\\

\verb! !\\
\verb! !\\




\end{multicols}
\end{document}